% Part: normal-modal-logic
% Chapter: syntax-and-semantics
% Section: normal-modal-logics

\documentclass[../../../include/open-logic-section]{subfiles}

\begin{document}

\olfileid{nml}{syn}{nor}

\olsection{نارمل امکانی منطقاں}

امکانی منطق نوں !!{formula}s دا کوئی ایہو جیہا سیٹ سمجھیا جا سکدا اے جیہڑا ڈھکدیاں شرطاں پوریاں کردا ہووے۔ کجھ امکانی منطقاں خاص دلچسپی رکھدیاں نیں۔ اوہناں دیاں تعبیراں تے ورتوں دلچسپ نیں، پر وڈی گل ایہ اے کہ نظریاتی طور اُتے اوہناں دی چنگی سمجھ اے۔ ایہناں نوں نارمل امکانی منطقاں آکھدے نیں۔

\begin{defn}\ollabel{defn:modal-logic}
اک \emph{امکانی منطق}~$\Log L$ بنیادی امکانی زبان دے !!{formula}s دا سیٹ اے جیہڑا
\begin{enumerate}
\item سارے قضیائی ہمیشہ سچے فارمولے رکھدا اے،
\item یکساں بدلی ہیٹھ بند اے، تے
\item اثباتِ مقدم ہیٹھ بند اے، یعنی جدوں وی $!A$ تے $(!A \lif
  !B) \in \Log L$، تاں $!B \in \Log L$ وی۔
\end{enumerate}
\end{defn}

\begin{ex}
!!{formula}s دے کجھ سیٹ، جیہڑے منّیا کہ بہتے دلچسپ نہیں، ایہ تعریف پوری کردے نیں تے ایس لئی امکانی منطقاں نیں:
\begin{enumerate}
\item سارے ہمیشہ سچے فارمولیاں دا سیٹ،
\item سارے !!{formula}s دا سیٹ (متضاد امکانی منطق)۔
\end{enumerate}
\end{ex}

\begin{defn}\label{defn:normal-modal-logic}
امکانی منطق~$\Log L$ اُدوں تے صرف اُدوں \emph{نارمل} اے جدوں
\begin{enumerate}
\item ایہدے وچ دوویں نیں:
\begin{align}
\tag{K} \Box(p \lif q) & \lif (\Box p \lif \Box q) \\
\tag{Dual} \Diamond p & \liff \lnot \Box \lnot p
\end{align}
\item تے ایہ لازمیت دے قاعدے ہیٹھ بند اے، یعنی جدوں وی $!A \in
  \Log L$، تاں $\Box !A \in \Log L$ وی۔
\end{enumerate}
\end{defn}

امکانی منطقاں دی تعریف دے دو مُکھ ڈھنگ سامنے آؤن گے۔ پہلا ثبوتی نظریے والا اے۔ ایس وچ اوہ !!{formula}s دا سیٹ ہوندی اے جیہڑے کجھ !!{derivation} نظاماں وچ اخذ کیتے جا سکدے نیں۔ دوجا نمونیاں دے نظریے والا اے۔ ایس وچ اوہ !!{formula}s دے سیٹ ہوندے نیں جیہڑے نمونیاں دیاں کجھ جماعتاں وچ معتبر نیں۔ ایہ عام قضیائی منطق (تے ایس گل لئی پہلے درجے دی منطق وی) والی صورت وانگ اے۔ ایتھے وی اسیں !!{formula}s دا سیٹ دو وکھرے ڈھنگاں نال وکھ کر سکدے آں، مثلاً اوہناں ہمیشہ سچے فارمولیاں دا سیٹ جیہڑے سچائی دیاں قدراں دی ہر قدر بندی وچ سچ نیں، تے اوہناں سارے !!{derivable} !!{formula}s دا سیٹ جیہڑے کسے !!{derivation} نظام وچ نیں۔ نارمل امکانی منطقاں نوں ایہناں دوہاں ڈھنگاں نال مقرر کرنا تعلقی نمونے تے مسلّماتی (ہلبرٹ طرز دے) ثبوتی نظام ورت کے ممکن اے۔
\end{document}
